\documentclass[11pt,a4paper]{article}
\usepackage[T1]{fontenc}
\usepackage[utf8]{inputenc}
\usepackage{lmodern,microtype,amsmath,amssymb,amsthm,booktabs,longtable,array}
\usepackage[margin=25mm]{geometry}
\usepackage[colorlinks=true,linkcolor=blue,citecolor=blue,urlcolor=blue]{hyperref}
\usepackage{url}
\usepackage[numbers]{natbib}
\hypersetup{pdftitle={Horizon : admission confidentielle sur deux engagements signes},pdfauthor={Dossier technique ZK-agentiK},pdfsubject={Demonstrations conditionnelles et verification numerique experimentale}}
\newtheorem{theorem}{Proposition}
\newtheorem{lemma}[theorem]{Lemme}
\newtheorem{definition}[theorem]{Définition}
\newtheorem{assumption}[theorem]{Hypothèse}
\newcommand{\Com}{\operatorname{Com}}
\newcommand{\Zq}{\mathbb Z_q}
\newcommand{\N}{\mathbb N}
\newcommand{\code}[1]{\nolinkurl{#1}}
\title{Horizon : admission confidentielle sur deux engagements signés\\
  \large Relation, preuves conditionnelles et limites expérimentales}
\author{Dossier technique du dépôt ZK-agentiK}
\date{17 octobre 2026 \quad Version expérimentale 1}
\begin{document}
\maketitle
\begin{abstract}
Horizon vérifie une inégalité entre un montant et un plafond par dossier,
attestés par deux rôles distincts, avant d'enregistrer une autorisation
synthétique. Il utilise Pedersen sur Ristretto, Bulletproofs et des signatures
Ed25519. Ce texte donne des démonstrations conditionnelles, sépare confidentialité,
authenticité et reprise, et fournit des vérifications arithmétiques exécutables.
Il ne constitue ni une preuve formelle machine, ni un audit cryptographique,
ni une nouvelle primitive, ni une qualification bancaire. Les essais Octave
ne prouvent pas la sécurité cryptographique. Le profil n'est pas post-quantique.
\end{abstract}
\tableofcontents
\newpage

\section{Objet, périmètre et correspondance avec le code}
Le profil \code{horizon-range-v1} se trouve dans
\code{labs/horizon/proof/src/lib.rs}. L'effet de démonstration est
\code{record_claim_authorization}, dans un registre SQLite local, pas un paiement.
La cible économique envisagée est une autorisation de sinistre déléguée ; aucune
demande solvable, conformité automatique ou supériorité sur Vega n'est établie.

La bibliothèque vérifie un \emph{énoncé cryptographique}. Le contrôleur
\code{packages/horizon-demo/src/store.ts} ajoute les attentes locales,
la consommation de nonce, la déduplication métier et la gestion d'état. Le SDK
de publication \code{packages/proofkit-sdk} reste un composant expérimental :
une interface ne remplace aucune de ces hypothèses.

\begin{center}\small
\begin{tabular}{p{.27\linewidth}p{.63\linewidth}}
\toprule
Objet & Statut de l'argument \\
\midrule
Identités algébriques et bornes & Démonstrations mathématiques ci-dessous. \\
Primitive et composition & Sécurité conditionnelle aux hypothèses explicites ; héritée des travaux cités. \\
Code et état local & Tests adverses et modèle fini abstrait ; pas de preuve de raffinement du code. \\
Octave/MATLAB & Vérification numérique exacte d'exemples et d'un domaine réduit ; exécution Octave distincte de MATLAB. \\
Manifestes SHA-256 & Intégrité relativement à un manifeste de confiance ; date auto-déclarée, pas horodatage tiers, signature ou preuve de sens. \\
\bottomrule
\end{tabular}\end{center}

\section{Paramètres et contrat d'émission}
Soit $\mathbb G$ le groupe Ristretto d'ordre premier
\begin{align*}
q&=2^{252}+27742317777372353535851937790883648493.
\end{align*}
Les encodages canoniques et le groupe premier sont spécifiés par RFC~9496
\cite{rfc9496}. $G,H\in\mathbb G$ sont les bases fixes de
\code{PedersenGens::default()}. Le prover ne les choisit pas. Posons $N=2^{32}$.

Deux autorités produisent $A=\Com(a,r_a)=aG+r_aH$ et
$C=\Com(c,r_c)=cG+r_cH$. Le contrat exige des valeurs entières non négatives
représentables en \code{u64}, la même devise EUR et le même exposant 2.
Les signatures couvrent profil, rôle, autorité, tenant, sujet, devise, exposant,
version, clé du détenteur, période et engagement. Les clés de confiance sont
épinglées ; les deux identifiants et clés d'autorité doivent être distincts.

\begin{assumption}[Émission et jointure]
Chaque autorité calcule ou contrôle l'ouverture de son engagement à partir
de ses faits. Elle ne signe pas aveuglément un point soumis par le détenteur.
L'autorité montant ne connaît pas nécessairement le plafond, et inversement.
Les deux ouvertures arrivent au prover ; celui-ci connaît les deux valeurs.
Les métadonnées communes désignent effectivement le même dossier et la même
version métier. Le contrôleur associe indépendamment ce sujet à l'objet ciblé.
\end{assumption}
Deux clés différentes ne prouvent pas l'indépendance organisationnelle, la
qualité de l'expertise ni l'absence de collusion. La version de schéma 1 du
prototype n'est pas un registre d'historique des sinistres. Une signature
atteste une émission sous clé, pas la vérité du fait signé.

\section{Relation et démonstrations conditionnelles}
\subsection{Engagement : confidentialité et liaison}
\begin{lemma}[Hiding dans le modèle idéal]
Pour $H\ne0$ et $r$ uniforme dans $\Zq$, la loi de $vG+rH$ est uniforme
dans $\mathbb G$ pour toute valeur fixée $v$.
\end{lemma}
\begin{proof}
La multiplication $r\mapsto rH$ est une bijection puisque $q$ est premier
et $H\ne0$. La translation par $vG$ conserve l'uniformité. Les distributions
des engagements sont donc égales pour deux valeurs quelconques.
\end{proof}
Ce résultat concerne l'engagement seul. Dans le code, la qualité du générateur
aléatoire et de la conversion en scalaires reste une hypothèse d'implémentation.
Il ne justifie pas des masques réutilisés ou choisis en fonction du secret.

\begin{lemma}[Liaison calculatoire]
Produire deux ouvertures $(v,r)\ne(v',r')$ du même engagement avec
$v\ne v'\pmod q$ permet de calculer $\log_G H$.
\end{lemma}
\begin{proof}
L'égalité donne $(v-v')G=(r'-r)H$. Si $r'=r$, alors $v=v'\pmod q$,
contradiction. Ainsi $(r'-r)^{-1}$ existe dans $\Zq$ et
$\log_G H=(v-v')(r'-r)^{-1}\pmod q$.
\end{proof}
La confidentialité est informationnelle dans le modèle idéal ; la liaison
dépend du logarithme discret et de bases sans relation connue
\cite{pedersen1992,daleknotes}. Aucune expérience numérique réduite ne valide
cette hypothèse pour Ristretto.

\subsection{Énoncé effectivement prouvé}
Les engagements dérivés sont
\[
 X=A-G,\qquad D=C-A.
\]
Le prover appelle une preuve de plages agrégée à deux valeurs de 32 bits
\cite{bulletproofs,dalek5}, avec ouvertures
\[
 x=a-1,\quad d=c-a,\quad r_X=r_a,\quad r_D=r_c-r_a\pmod q.
\]
La relation d'ouverture est
\begin{equation}\label{eq:relation}
\exists x,d\in\{0,\ldots,N-1\},\ r_X,r_D\in\Zq:
 X=xG+r_XH,\quad D=dG+r_DH.
\end{equation}
\textbf{Une nuance nécessaire :} comme $H$ engendre le groupe, tout point
possède mathématiquement une ouverture à n'importe quelle valeur. Une simple
affirmation d'appartenance existentielle ne suffit donc pas à lier une valeur
métier. Le raisonnement utilise l'\emph{argument de connaissance} d'ouvertures
et la liaison calculatoire aux ouvertures effectivement attestées.

\begin{assumption}[Primitive héritée]
Le protocole Bulletproofs non interactif utilisé fournit la complétude,
la connaissance calculatoire des ouvertures et la confidentialité appropriées
à son modèle de sécurité et à ses transcriptions Fiat--Shamir.
Son implémentation, ses paramètres, son aléa et son analyse de transcript
correspondent à ce protocole. La version corrigée du papier est la référence ;
ce dossier ne refait ni sa réduction ni un audit de la bibliothèque.
\end{assumption}
Le papier ePrint signale une correction de sa section Fiat--Shamir en 2022
\cite{bulletproofs}. Cette histoire interdit d'assimiler une citation de
sécurité à une preuve automatique de toute adaptation logicielle.

\begin{theorem}[Complétude conditionnelle]
Si les signatures et métadonnées sont valides, la politique permet le mode,
le challenge est celui attendu et non expiré, et
$1\le a\le N$, $a\le c<a+N$, alors un prover honnête construit les deux
engagements dérivés avec des valeurs admissibles. Sous la complétude de la
primitive et en l'absence d'erreur d'exécution, sa preuve est acceptée.
\end{theorem}
\begin{proof}
$x=a-1\in[0,N-1]$ et $d=c-a\in[0,N-1]$. Par linéarité,
$X=(a-1)G+r_aH$ et $D=(c-a)G+(r_c-r_a)H$.
Ce sont exactement les entrées attendues par le vérifieur. Les autres
préconditions sont testées séparément par celui-ci.
\end{proof}
Cette complétude ne garantit pas disponibilité réseau, absence de révocation
après calcul, ni succès d'un effet métier.

\begin{theorem}[Bornes entières et absence de wrap]\label{th:bounds}
Supposons l'extraction des ouvertures de \eqref{eq:relation}, la liaison
des engagements et l'émission dans le domaine entier \code{u64}. Alors les
valeurs attestées satisfont
\[
1\le a\le2^{32},\qquad 0\le c-a<2^{32},\qquad c\le2^{33}-1.
\]
\end{theorem}
\begin{proof}
Les ouvertures extraites donnent une ouverture de $A$ à $a'=x+1$ avec
masque $r_X$, et de $C$ à $c'=x+1+d$ avec masque $r_X+r_D$.
Elles satisfont $1\le a'\le N$ et $a'\le c'\le2N-1$.
Si $a'$ ou $c'$ diffèrent modulo $q$ des valeurs attestées, on obtient une
violation de liaison. Sinon $a=a'\pmod q$ et $c=c'\pmod q$.
Les valeurs attestées sont entre 0 et $2^{64}-1<q$, et les extraites entre
0 et $2^{33}-1<q$. Deux entiers de $[0,q)$ congrus sont égaux ; les bornes
sont donc des bornes entières, et non seulement modulaires.
\end{proof}
Le plafond peut dépasser $2^{32}-1$ : il faut \code{u64}. Le maximum
$8\,589\,934\,591$ centimes vaut $85\,899\,345{,}91$ EUR, sans engagement
sur l'adéquation de ce maximum à un contrat d'assurance.

\begin{theorem}[Correction calculatoire conditionnelle]
Conditionnellement aux autorités honnêtes pour l'émission, à la liaison,
à l'argument de connaissance hérité, à l'inforgeabilité des signatures
et à une implémentation fidèle, accepter une valeur attestée qui viole les
bornes de la proposition~\ref{th:bounds} exige de rompre au moins l'une de
ces hypothèses ou les vérifications de liaison des contextes.
\end{theorem}
\begin{proof}
Une signature sur une attestation non émise relève d'une contrefaçon.
Pour les attestations émises, l'extraction et la proposition précédente
imposent les bornes. Une mauvaise association de contexte exige soit une
entrée de confiance erronée, soit une erreur d'encodage/comparaison,
soit une rupture du mécanisme qui authentifie ce contexte.
\end{proof}
On peut organiser une borne d'échec comme une somme des événements de
contrefaçon, de collision, d'échec d'extraction et de rupture de liaison.
Sans réduction complète et modèle des oracles, ce dossier ne donne ni
constantes de réduction ni probabilité numérique de sécurité. Les 500
rejets d'une campagne ne sont pas une estimation de la soundness.

\section{Composition avec signatures et challenge}
Le challenge attendu est émis et stocké par le contrôleur. Il contient tenant,
client, audience, action, opération, objet, sujet, politique et époque, époque
de confiance, nonce de 32 octets, dates et méthode. L'égalité porte sur
\emph{tous} ces champs, avant autorisation. L'agent ne fournit pas les racines
de confiance. Les encodages sont du JSON typé en ordre de déclaration,
avec domaines distincts ; ce n'est pas une promesse de JSON canonique universel.

\begin{theorem}[Liaison de la présentation]
Supposons Ed25519 inforgeable pour les clés honnêtes, SHA-256 résistant aux
collisions, les domaines distincts et l'encodage injectif sur les types
admis. Une présentation acceptée sous un challenge indépendant est liée
au couple d'attestations, au payload et à ce challenge par la clé du détenteur.
\end{theorem}
\begin{proof}
La signature du détenteur porte sur le digest du tuple profil, attestations,
challenge et payload, lui-même encodé sous un domaine distinct. Changer un
de ces objets sans obtenir une nouvelle signature exige une collision ou
une contrefaçon. Une entrée dotée d'un autre challenge est rejetée par
l'égalité avec la valeur attendue. Les signatures des autorités couvrent
chaque corps d'attestation séparément.
\end{proof}
RFC~8032 décrit EdDSA \cite{rfc8032} ; le code utilise une vérification stricte
et refuse les clés faibles. Une clé compromise ou une autorité malveillante
sort de l'hypothèse, et deux présentations identiques peuvent toujours être
rejouées en l'absence d'un état de consommation.

\paragraph{Mode A.} Le challenge entier est aussi ajouté au transcript
Merlin. La preuve de plages est recalculée par présentation. Il s'agit d'une
liaison cryptographique au contexte, pas d'une horloge ou d'une preuve de
vérité instantanée des sources.
\paragraph{Mode B.} La preuve de plages est préparée sur les attestations
immuables puis conservée. Une nouvelle signature du détenteur couvre le
challenge et la preuve exacte. Cela atteste une autorisation fraîche par
cette clé sur une preuve historique, \emph{pas} une connaissance fraîche du
witness. Le vérifieur actuel vérifie toute la preuve à chaque appel.
\paragraph{Politique C.} Si la confidentialité n'est pas exigée, le sélecteur
choisit la divulgation signée. Sinon, il choisit B si la réutilisation est
autorisée et A dans le cas contraire. Un refus cryptographique n'entraîne
aucun fallback. L'autorisation des méthodes reste vérifiée indépendamment.

\section{Anti-rejeu et frontière transactionnelle}
L'identité d'opération du contrôleur est le hachage du tuple tenant, audience,
action et objet. Elle ne contient pas le nonce. La garantie est locale à ce
domaine de déduplication ; plusieurs audiences ou bases indépendantes ne
partagent pas automatiquement une identité métier globale.

\begin{assumption}[Modèle de stockage et de cible]
Les transactions sont sérialisées et atomiques. Les états commis sont
persistés sans rollback de sauvegarde, effacement administratif ou corruption.
Tous les chemins d'effet passent par le même contrôleur, la même identité
métier stable et la même base. Le résultat \code{Verified} provient du worker
cryptographique de confiance, jamais d'un champ HTTP forgé. L'horloge et
les racines de confiance sont contrôlées par le vérifieur.
\end{assumption}

\begin{theorem}[Au plus un effet local par identité métier]
Sous l'hypothèse précédente, le contrôleur produit au plus une ligne d'effet
pour une identité d'opération, y compris avec plusieurs nonces, concurrence,
arrêt après réservation ou réponse ambiguë.
\end{theorem}
\begin{proof}
Initialement, aucune opération n'existe. Une transaction valide le contexte
courant, consomme le nonce et insère l'opération sous une clé primaire unique.
Toute réservation ultérieure avec cette clé est rejetée. La transaction
de cible exige l'état réservé, recontrôle échéances et époques, puis insère
l'effet sous la même clé primaire et termine l'opération. Un échec de ce
recontrôle conserve une opération rejetée ; une réponse ambiguë conserve
un état inconnu. Aucun de ces états n'est supprimé pour retenter. La
réconciliation consulte les effets existants et ne redispatche pas.
Par induction sur les transitions, le nombre d'effets reste au plus un.
\end{proof}
SQLite fournit les mécanismes transactionnels sous ses hypothèses documentées
\cite{sqlite,sqliteatomic}. La preuve est un invariant abstrait, pas un
théorème vérifié sur le moteur, le filesystem ou le code TypeScript.
Elle ne prouve ni qu'un effet finira par arriver, ni une atomicité sur une
API distante. Un arrêt après réservation peut laisser une opération bloquée.
Le journal chaîné détecte certaines modifications si une racine est connue,
mais un administrateur capable de réécrire toute la base n'est pas exclu
par un simple hachage local.

La validité temporelle est l'intersection des intervalles :
\[
t_{\min}=\max(t_{\mathrm{challenge}},t_A,t_C),\qquad
t_{\max}=\min(e_{\mathrm{challenge}},e_A,e_C),\qquad
t_{\min}\le t<t_{\max}.
\]
Les revocations et époques sont recontrôlées après l'attente cryptographique
et avant l'effet. Cela borne la validité d'un snapshot, sans prouver la
fraîcheur d'un système externe indisponible en crise.

\section{Fuites et contre-exemples métier}
Les valeurs et ouvertures sont privées face au vérifieur dans A/B, sous les
hypothèses de la preuve. Les autorités connaissent leurs propres valeurs,
le prover connaît les deux, et le mode divulgué les révèle toutes.
Engagements, signatures, clés, dossier opaque, tenant, rôles, dates, politique
et décision sont publics. Des engagements fixes rendent les sessions
corrélables ; rerandomiser la preuve seule ne les rend pas anonymes.
La décision apprend au moins l'inégalité et les bornes. Des politiques
adaptatives, informations auxiliaires et canaux temporels peuvent révéler
davantage ; aucune non-interférence système n'est démontrée.

\begin{center}\small
\begin{tabular}{p{.19\linewidth}p{.35\linewidth}p{.36\linewidth}}
\toprule
Contre-exemple & Énoncé accepté & Ce qui manque \\
\midrule
Budget cumulé & Deux dossiers à 60 satisfont chacun $60\le100$. & Leur somme 120 excède un budget commun de 100 ; il faut un ledger concurrent. \\
Franchise & Un versement 15000 est inférieur à 20000. & Si perte 15000 et franchise 5000, le montant prévu serait 10000 ; cette formule n'est pas prouvée. \\
Exposant & Les entiers $1000\le10000$ passent. & 1000 EUR dépasse 10000 centimes ; le prototype empêche ce mélange en fixant EUR/exposant 2. \\
Devise & Deux nombres égaux. & EUR et USD ne sont pas une unité commune ; conversion, source et instant doivent être authentifiés. \\
Wrap & $103\equiv2\pmod{101}$. & Un domaine source non borné peut changer le sens entier ; le contrat \code{u64} exclut ce cas pour le vrai $q$. \\
Nonce seul & Deux nonces différents sont chacun neufs. & Ils peuvent autoriser deux fois le même objet sans réservation métier. \\
\bottomrule
\end{tabular}\end{center}
Ces exemples ne sont pas des contournements des contrôles présents ; ils
délimitent l'énoncé. Une logique de sinistre complète doit expliciter exclusions,
franchises, taxes, coassurance, dates de couverture et réserves. Aucun circuit
ne devient « parfait » en augmentant seulement sa vitesse.

\section{Mesures CPU, GPU, Vega et hypothèses quantiques}
Les mesures conservées dans \code{docs/horizon/BENCHMARK_RESULTS.json}
utilisent 500 observations après 30 échauffements. Les fonctions crypto Linux
parcourent cinq cas ; les quantiles sont au rang supérieur. L'indépendance IID
des observations n'est pas établie et les IC lui sont conditionnels.
Les mêmes paramètres publics fixes sont initialisés une fois par processus
avec \code{OnceLock} : ni preuve, ni witness, ni verdict, ni aléa n'y sont cachés.
Le mode B conserve séparément sa preuve préparée, conformément à sa sémantique.

\begin{center}\small
\begin{tabular}{p{.46\linewidth}p{.4\linewidth}}
\toprule
Mesure publiée du laboratoire & Observation, pas SLA \\
\midrule
A, génération CPU p95 avant / après & 7,770 / 6,546 ms \\
B, vérification complète CPU p95 avant / après & 2,266 / 1,247 ms \\
B, présentation signée p95 après préparation & 0,065 ms ; pas une nouvelle preuve \\
Preuve agrégée seule & 672 octets ; enveloppe et signatures exclues \\
Baseline divulguée, vérification p95 & 0,245 ms \\
Windows IPC, vérification et effet local durable & p95 13,014 ms ; maximum 180,310 ms, hors HTTP \\
Vega MC SHA-256 2 Kio, 8 threads & Preuve 65,836 ms, vérification 19,885 ms ; estimations centrales Criterion, pas p95 \\
\bottomrule
\end{tabular}\end{center}
La taille de preuve est ici $32(2\log_2(2\cdot32)+9)=672$ octets pour
le format agrégé utilisé. Elle n'inclut pas les engagements publics ou les
autorisations. La réduction observée du p95 par conservation des paramètres
est 15,75\% pour A et 44,97\% pour la vérification B ; aucun record externe
n'en découle.

Vega s'intéresse à des présentations de credentials et à une préparation
réutilisable \cite{vega,vegacode}. Le bench SHA-256 du dépôt public n'est pas
le circuit applicatif complet des 92/23 ms du papier. Horizon exige au contraire
une émission adaptée sur des engagements et laisse signatures et métadonnées
publiques. Les relations, modèles de fuite et phases diffèrent : aucun rapport
de vitesse ne permet de conclure que l'un bat l'autre sur une tâche commune.

Le dossier GPU est conservé séparément dans \code{docs/experiments/gpu/}.
Une tentative CUDA/RISC Zero n'est pas un portage GPU de cette relation
Bulletproofs. Ce texte n'attribue aucune accélération GPU sans résultat
reproduit, déclaration de relation, taille, transfert hôte/appareil, préparation,
vérification indépendante et matériel exact. Une indisponibilité CUDA est
un résultat de faisabilité négatif, pas une mesure de vitesse. Les chiffres
GPU ne sont donc pas interpolés depuis ceux CPU.

Enfin, GPU et résistance quantique sont des axes différents. Un calcul de
logarithme discret efficace compromettrait la liaison et les signatures
elliptiques ; l'algorithme de Shor motive cette distinction \cite{shor}.
ML-DSA et SLH-DSA sont des standards de signature post-quantique
\cite{fips204,fips205}. Remplacer Ed25519 par l'un d'eux ne rendrait pas
Pedersen/Bulletproofs post-quantiques. Le fait d'utiliser SHA-256 ou SHA-3
dans un transcript ne transforme pas un protocole elliptique en preuve
entièrement fondée sur des hachages. Une migration demande une autre relation
de preuve, une analyse de composition et de nouveaux benchmarks.

\section{Vérifications reproductibles et limites de preuve}
\code{math/horizon_checks.m} énumère les paires entières d'un domaine réduit
pour 1 à 6 bits, incluant des valeurs négatives ; contrôle les frontières
32 bits en \code{uint64} ; reconstruit $q$ par chaînes décimales ; et explore
tous les états atteignables d'un modèle fini de réservation à deux nonces
et une opération. Les fonctions de calcul décimal ne manipulent que de petits
chiffres, jamais $q$ en flottant. Aucune implémentation de courbe, signature
ou Bulletproofs n'est recréée dans MATLAB/Octave.

Le résultat JSON indique explicitement runtime, version, date UTC et limites.
Le conteneur fige le digest de base et le dossier conserve l'identifiant de
l'image construite et les versions des paquets. Le Dockerfile dépend encore
du dépôt de paquets au moment de sa reconstruction : ce n'est pas une
promesse de build reproductible bit à bit. L'exécution est sans réseau après
construction. MATLAB nécessite un runtime/licence distinct : une réussite
Octave n'est pas présentée comme une exécution MATLAB.

Les tests adverses Rust et Node existants restent une preuve expérimentale
séparée. Le modèle réduit ne remplace ni la vérification d'une vraie preuve
cryptographique par un processus indépendant, ni un modèle-checker appliqué
au code, ni une revue de sécurité indépendante. Le journal de compilation PDF
et l'inspection des pages attestent seulement la fabrication du document.

\section{Décision technique}
Le résultat défendable est un laboratoire de faisabilité : une relation étroite,
des hypothèses explicites, des preuves réelles et un contrôleur local qui
refuse avant son effet. La sécurité de production exige encore autorités
qualifiées, gestion des clés, revocation, isolation, audit, modèle distant de
reprise et validation du besoin économique. Les déploiements cloud restent
des exemples de configuration réseau/IAM, pas une certification opérationnelle.
Une affirmation de conformité, d'absence de concurrence ou de logique métier
universellement correcte ne découle d'aucun théorème présenté ici.

\bibliographystyle{plainnat}
\bibliography{references}
\end{document}
